\section{Mecanismos de defensa}

\subsection{Principios de defensa}

\begin{importantbox}{Tres principios de defensa}
\begin{enumerate}
    \item \textbf{Defensa en profundidad} (Defense in Depth): el modelo del queso suizo, es decir, defensa en múltiples capas, de modo que aunque una capa falle, las demás capas aún pueden bloquear el ataque. Las capas incluyen: depuración de entradas $\to$ endurecimiento del modelo $\to$ aplicación de políticas de acciones $\to$ monitoreo y detección de anomalías
    \item \textbf{Privilegio mínimo y separación de privilegios} (Least Privilege \& Privilege Separation): el sistema solo debe contar con el privilegio mínimo necesario para cumplir su función
    \item \textbf{Seguridad desde el diseño} (Safe/Secure by Design): en el caso ideal, la verificación formal ofrece garantías de seguridad demostrables
\end{enumerate}
\end{importantbox}

\subsection{Mecanismos de defensa concretos}

\subsubsection{Endurecimiento del modelo}

El endurecimiento se aplica en cada etapa del desarrollo del modelo:
\begin{itemize}
    \item Etapa de depuración y preparación de datos
    \item Etapa de preentrenamiento seguro
    \item Etapa de alineación posterior al entrenamiento (Post-training Alignment)
    \item Desaprendizaje automático (Machine Unlearning)
\end{itemize}

\subsubsection{Barreras de depuración de entradas}

Se revisan las entradas del LLM: se comparan con criterios predefinidos, se escapan los caracteres especiales y se normalizan a un formato estándar, con el fin de bloquear entradas maliciosas.

\subsubsection{Aplicación de políticas de acciones}

\begin{importantbox}{Marco de control de permisos programable}
El equipo de Dawn Song desarrolló un marco de control de permisos programable para agentes LLM:
\begin{itemize}
    \item Se revisa cada llamada a herramienta del agente conforme a una política
    \item Por ejemplo: se prohíbe la operación \texttt{delete\_db} y solo se permite cargar bases de datos específicas
    \item Aunque la base de conocimiento RAG sea envenenada y lleve al agente a intentar eliminar la base de datos, la capa de aplicación de políticas bloquea esa operación
\end{itemize}
\end{importantbox}

El caso de una aplicación bancaria muestra cómo, mediante políticas, se limita el monto máximo de transferencia y el rango de destinatarios, para impedir que el agente sea manipulado y ejecute operaciones financieras fuera de lo autorizado.

\subsubsection{Barreras de salida}

Se revisan las salidas del LLM para impedir que contenido dañino, tóxico o inapropiado llegue al usuario o a sistemas externos.

\subsubsection{Monitoreo y detección de anomalías}

Se monitorea de manera continua el comportamiento del agente, se detectan patrones anómalos y se identifican y responden a tiempo posibles ataques.

\subsection{Resumen de la sección}

La defensa de la IA agéntica requiere una estrategia integral de múltiples niveles y múltiples mecanismos que, desde el endurecimiento del modelo hasta la depuración de entradas, el control de acciones, el filtrado de salidas y el monitoreo continuo, conforme un sistema completo de defensa en profundidad.

